% =====================================================================
% 附录 B：实验设置与构造细节。由原 appendix.tex 的 A.3/A.4 搬运而来。
% =====================================================================
\section{Setup and construction details}
\label{app:method}

\subsection{Complete experimental setup}
\label{app:setup-details}
\paragraph{Scope and shared settings.} We evaluate twelve dataset--backbone settings over nine public datasets:
MNIST, ImageNet-100, Cora, IMDb, AG News, California Housing, STS-B,
ZINC-12k, and AgeDB, covering MLP, CNN, GCN, and RNN backbones. All
methods use the same data splits ($60/20/20$ train/validation/test
with a fixed seed for most datasets, official splits kept where provided,
e.g., Cora and ZINC-12k), backbone, Gaussian posterior with learned diagonal
variance, bottleneck $d=256$, Adam training, early stopping, and
validation-based selection; GPB's prior covariance is the
fixed shared $\tau^2I$ ($\tau^2=1$).

\paragraph{Training defaults and preprocessing.} The shared training defaults are
hidden sizes $512$--$256$, dropout $0.2$, learning rate $10^{-3}$ (Cora
uses $10^{-2}$), batch size $512$, and a $200$-epoch budget with early
stopping (patience $20$); dataset-specific defaults are likewise shared
by every method on that dataset: STS-B uses a single-layer LSTM with
hidden size $256$ and frozen GloVe embeddings; AG News uses
$50$-dimensional PCA-reduced BERT features; ImageNet-100 uses $4\times4$
spatial pool-4 ResNet-50 features; the remaining dataset-specific
preprocessing is given in the release. Every configuration is trained
on the same ten seeds ($0$--$9$). Input features are standardized per
dataset, and regression targets are min--max normalized to $[0,1]$;
both are fitted on the training split only, and reported MAE is
rescaled to the original units.

\paragraph{Baselines and search-budget fairness.} The baselines are Base, VIB, SVIB,
NIB, CEB, supervised $\beta$-DVCCA, and the external non-IB baselines
NCBD and AdaCap: NCBD trains unit-hypersphere features and prototypes
with the NONL contrastive objective (simplex-ETF optimum), and AdaCap
pairs a permutation contrastive loss with a Tikhonov closed-form readout,
so the comparison spans geometric/prototype and contrastive non-IB
mechanisms in addition to the IB variants. GPB is OPB on classification
rows and EPB on regression rows. The main
comparison equalizes the search budget: the plain backbone has no
method-specific hyperparameters (one configuration); every
$\beta$-weighted method, VIB, SVIB, NIB, CEB, and DVCCA, is swept over
the same 30-point $\beta$ grid (six decades, five values per decade);
GPB searches the two-dimensional grid
$\beta\in\{10^{-4},10^{-3},10^{-2},10^{-1},1,10\}\times
a=\rho\in\{1,3,6,9,12\}$, also $6\times5=30$ configurations; NCBD and
AdaCap use 30-value grids of their own. The guarantees of Section 2.2
are stated for any fixed scale value and apply verbatim to the selected
configuration, and the same validation-based selection rule applies to
every ablation variant (Appendix~\ref{app:ablation}). The per-method
grids are summarized in Table~\ref{tab:search-protocol}.

\begin{table}[t]
\centering\small
\setlength{\tabcolsep}{3pt}
\caption{Per-method hyperparameter grids of the main comparison
(Table~\ref{tab:main_result}). Each configuration is trained on the same
ten seeds under the shared training budget; the configuration with the
best mean validation metric is selected per method and setting. Every
method evaluates the same search budget of thirty configurations:
$\beta$-weighted methods sweep a 30-point $\beta$ grid, GPB sweeps the
$\beta\times a=\rho$ grid below, and NCBD and AdaCap use 30-value grids
of their own.}
\label{tab:search-protocol}
\begin{tabular}{l l c}
\hline
Method & Method-specific hyperparameters & Configs. \\
\hline
Base & --- & $1$ \\
VIB, SVIB, NIB, CEB, DVCCA & $\beta\in\{1,2,4,6,8\}\times10^{k}$, $k=-4,\dots,1$ & $30$ \\
NCBD & temperature $T$: 30 log-spaced values in $[0.05,1]$ & $30$ \\
AdaCap & Tikhonov $\lambda\in\{1,2,4,6,8\}\times10^{k}$, $k=-2,\dots,3$ & $30$ \\
GPB (OPB / EPB) & $\beta$ grid above $\times$ $a=\rho\in\{1,3,6,9,12\}$ & $6\times5=30$ \\
\hline
\end{tabular}
\end{table}
\subsection{Construction details}
\label{app:construction}

\paragraph{Initialization, monitoring, and re-initialization.} The construction is
equivariant under label permutations ($Q_\theta(U_\theta P)=Q_\theta(U_\theta)P$),
so no class index plays a distinguished role; the posterior logvar
head is zero-initialized ($\sigma^2 = \tau^2 = 1$ at initialization, so training
starts at $\KL \approx \tfrac{1}{2\tau^2} a^2$); the prior encoder is
identity-initialized, so the initial raw matrix is $U_\theta = [e_1, \dots, e_K]$,
full column rank, with a stable polar-factor backward (the verified full-rank
condition). The full-rank condition is \emph{verified, not
imposed}: $\sigma_{\min}(U_\theta)$ is monitored along training and stays bounded
away from zero in the reported runs, and no additional regularizer is added to the
objective, which remains Eq.~\eqref{eq:pgib-objective} exactly. Near rank deficiency
the polar factor's gradient is delicate; should the monitor flag a near-deficient
step, the recommended handling is a re-initialization of the prior encoder (the
reported runs never triggered it).

\paragraph{Covariance protocol.} Equations~\eqref{eq:opb-prior} and
\eqref{eq:epb-prior} define the OPB and EPB instances used throughout the paper.
Their prior covariance is the fixed shared $\tau^2I$ with $\tau^2=1$ (matching the
zero-initialized posterior at $\sigma^2=1$ so training starts at
$\KL\approx a^2/(2\tau^2)$), with no learned prior-covariance parameters. The
posterior retains its learned diagonal covariance $\sigma^2_\phi(x)$, as specified
in Eq.~\eqref{eq:pgib-covariance}. The anchor-energy readout is the exact Gaussian
Bayes head of the benchmark model: the scores of Eq.~\eqref{eq:opb-energy} are
Gaussian Bayes scores up to a class-independent term, and for uniform classes the
corresponding linear weights are $(a/\tau^2)q_{\theta,k}$. The direct diagnostics
in Section~\ref{sec:direct-mismatch}
normalize by the declared $\tau^2$.

\paragraph{EPB axis choices.} Three variants of the direction: (A) \emph{fixed
axis}, $u = [1, 0, \dots, 0]$, with no trainable prior direction: the
theoretically cleanest ablation, at the price of prespecifying a latent coordinate;
(B) \emph{trainable normalized axis}, $u = W/\|W\|$ with $W$ trained and normalized
at every step, our main variant, which keeps the scale $\rho$ exact while allowing
the model to choose the axis; (C) \emph{unconstrained linear axis},
$m_\theta(y) = W \tilde y$: the simplest ablation, still isometric by linearity but
with a free scale that can shrink and weaken the geometry. We recommend (B); the
frozen-prior variant EPB-RV (Appendix~\ref{app:ablation},
Table~\ref{tab:ablation-ci-epb}) shows the trainable rotation carries no benefit
on Housing, and we keep (C) as a diagnostic. The main variant uses no bias: an affine offset cancels in pairwise
distances but is a free extra degree of matching, and $\tilde y$ is already centered.

\paragraph{EPB isometry notes.} The isometry is exact and global: any nonzero linear
map of a scalar label is isometric up to its norm, and the normalization fixes that
norm to $\rho$; unlike fixed random-feature anchors, whose pairwise distances grow
nonlinearly and saturate, the isometric axis preserves label distances at every scale
(with the price that anchor norms grow with $|\tilde y|$; a line and a fixed-radius
sphere intersect in at most two points, so isometry and equal norms cannot both hold).
All prior means lie on one line; for $y_i < y_j < y_k$ the distances add,
$\|m_\theta(y_i)-m_\theta(y_k)\| = \|m_\theta(y_i)-m_\theta(y_j)\| + \|m_\theta(y_j)-m_\theta(y_k)\|$:
order and magnitude of label differences are encoded exactly.

\paragraph{Multidimensional regression labels.} For $y \in \R^m$ ($m \le d$), the
same construction extends verbatim: take $W \in \R^{d \times m}$, its thin polar
factor $Q_\theta = W (W^{\top}W)^{-1/2}$, and the prior $m_\theta(y) = \rho\, Q_\theta\, \tilde y$,
which is isometric in the label metric; the scalar case is $m = 1$ with $Q_\theta = u$. The
label metric must be declared before standardization, otherwise ``isometric'' has no
unique meaning.
